%%%%%%%%%%%%%%%%%%%%%%%%%%%%%%%%%%%%%%%%%%%%%%%%%%%%%%%%%%%%%%
%%%%%%%%%%%% MA-0515 Ecuaciones Diferenciales %%%%%%%%%%%%%%%%
%%%%%%%%%%%%%%%%%%%%%%%%%%%%%%%%%%%%%%%%%%%%%%%%%%%%%%%%%%%%%%
\newpage
\subsection*{MA-0515 Ecuaciones Diferenciales}
\addcontentsline{toc}{subsection}{MA-0515 Ecuaciones Diferenciales}

\begin{refsection}


\begin{table}[H]
\centering{
\begin{tabular}{|ll|}
\hline
%\multicolumn{2}{|c|}{\textbf{Nivel de virtualidad del curso:} Virtual}\\ \hline
\textbf{Año:} III  & \textbf{Requisitos:} MA-0471 y MA-0515 \\ 
\textbf{Ciclo:} VI  & \textbf{Correquisitos:} MA-0635 \\ 
\textbf{Tipo de curso:} Teórico & \textbf{Horas presenciales por semana:}   \\ 
\textbf{Créditos:} 4 & \textbf{Horas de trabajo independiente por semana:}  \\ \hline
\end{tabular}
}
\end{table} 


\subsubsection*{Descripción del curso}

En este curso, que se ubica en el sexto ciclo del plan de estudios, proporciona los fundamentos teóricos de las ecuaciones diferenciales ordinarias. Estas estructuras matemáticas, que son ubicuas en aplicaciones, ya fueron discutidas, desde un punto de vista más intuitivo y operacional, en el curso ... El presente curso debería proveer el sustento necesario para dar lugar a estudios más avanzados en el tema así como a aplicaciones más complejas de la teoría. 

La teoría de existencia y unicidad fue desarrollada en el siglo XIX, mucho después de la introducción de las ecuaciones diferenciales en la matemática. Sus principales propulsores fueron Cauchy, Lindelöf, Lipschitz y Picard. El que podría considerarse como el principal teorema del curso es un elegante resultado que proporciona un conjunto de condiciones suficientes para garantizar existencia, unicidad y continuidad con respecto a parámetros en las ecuaciones diferenciales ordinarias. Posteriormente, pero en la misma línea, el tema de estabilidad le proporcionará a la persona estudiante una introducción a los sistemas dinámicos, la cual es otra área fascinante de la matemática y llena de aplicaciones.

Una segunda parte del curso con un sabor bastante distinto es la de problemas de valor de frontera (teoría de Sturm-Liouville). Esta teoría, aunque se desarrolla en el caso unidimensional, introducirá a la persona estudiante a métodos que son más bien característicos de ecuaciones en derivadas parciales. De hecho, aunque el objeto de estudio sigue siendo el de ecuaciones diferenciales ordinarias, aquí se encontrarán más similitudes con la teoría espectral del álgebra lineal que con los sistemas dinámicos de la primera parte del curso.

Para el buen desempeño de la persona estudiante en este curso, se espera que tenga dominio en los temas básicos de cálculo y análisis real en una y varias variables. También serán relevantes los conocimientos previos en álgebra lineal (en especial la teoría de valores y vectores propios).



\subsubsection*{Objetivos}

\textbf{General:}Establecer criterios para determinar la existencia, unicidad y estabilidad de soluciones de Ecuaciones Diferenciales Ordinarias\\

\textbf{Específicos:}

\begin{enumerate}
\item Identificar bajo cuales condiciones se tiene existencia, unicidad y estabilidad de las soluciones.
\item Establecer la existencia de soluciones maximales y de tipos de dependencia de las soluciones respecto a las condiciones iniciales.
\item Determinar la estabilidad de un punto de equilibrio de un sistema de ecuaciones diferenciales.
\item Analizar un problema de valor de frontera para determinar la existencia de soluciones y las propiedades de las mismas.
\end{enumerate} 

\subsubsection*{Contenidos}


%Teoremas de existencia y unicidad (Picard y Peano), soluciones maximales y comportamiento en la frontera, dependencia de valores iniciales y parámetros, estabilidad de Liapunov, problemas de valores de frontera, funciones de Green, problemas de Sturm-Liouville

\begin{enumerate}


\item \textbf{Introducción y ecuaciones diferenciales ordinarias de primer orden  (2 semanas):} 
\begin{enumerate}
    \item Definiciones y terminología: clasificación por tipo (ordinaria o parcial), clasificación por orden, clasificación por linealidad. Definición de solución de una EDO, soluciones explícitas e implícitas.
    \item Problemas de valor inicial. Definiciones de existencia y unicidad de soluciones.
    \item Ecuaciones de variables separables.
    \item Ecuaciones lineales.
    \item Ecuaciones exactas.
    \item Solución de ecuaciones por sustitución.
\end{enumerate}

\item \textbf{Sistemas de ecuaciones diferenciales lineales (3.5 semanas):} 
\begin{enumerate}
    \item Repaso de valores propios.
    \item Sistemas lineales homogéneos: valores propios reales distintos, valores propios repetidos y valores propios complejos.
    \item Solución mediante diagonalización.
    \item Exponencial de una matriz y valores propios repetidos
    \item Sistemas lineales no homogéneos: matriz fundamental, coeficientes indeterminados, variación de parámetros y diagonalización.
    \item Ecuaciones diferenciales lineales ordinarias de orden superior
\end{enumerate}
\item \textbf{La transformada de Laplace (2.5 semanas):}
\begin{enumerate}
    \item Definición.
    \item Transformadas inversas y transformadas de derivadas.
    \item Propiedades de traslación.
    \item Derivada de una transformada, transformadas de integrales, convolución y transformadas de funciones periódicas.
    \item La función Delta de Dirac.
    \item Propiedades relativas a la invertibilidad
    \item Aplicación de la transformada de Laplace a la solución de sistemas de ecuaciones diferenciales e integro-diferenciales. 
    
\end{enumerate}
\item \textbf{Teoría de existencia y unicidad (2 semanas):} 
\begin{enumerate}
    \item Teoremas de existencia para ecuaciones en variables separables. 
    \item Teorema de Picard y aproximaciones sucesivas.
    \item Teorema de Peano.
\end{enumerate}


\item \textbf{Continuación de soluciones (1 semana):} 
\begin{enumerate}
    \item Solución Maximal.
    \item Comportamiento en la frontera.
\end{enumerate}

\item \textbf{Dependencia de valores iniciales y parámetros (2 semanas):} 
\begin{enumerate}
    \item Dependencia continua (local y global).
    \item Dependencia diferenciable (local y global).
\end{enumerate}


\item \textbf{Tema a elegir por la persona docente (2 semanas):}
 Se dejan estas tres semanas para que la persona docente pueda introducir algún tema complementario al curso. Por ejemplo, algunas opciones son: Estabilidad, problemas de valor de frontera, modelación, ecuaciones en diferencias, ecuaciones en derivadas parciales, etc.


\end{enumerate}

\subsubsection*{Metodología}

\noindent \textbf{Lineamientos metodológicos:} La persona docente a cargo de este curso debe respetar las siguientes pautas:
\begin{enumerate}
    \item Fomentar el desarrollo de intuición sobre los conceptos estudiados en el curso. 
    \item Requerir de parte del estudiantado, una comunicación clara y rigurosa de argumentos y pruebas matemáticas de manera oral y escrita.
    \item Cuando la naturaleza del contenido lo permita, se deben utilizar representaciones visuales que apoyen la comprensión de los temas.
    \item Fomentar una visión integral de las matemáticas, específicamente discutir cómo se aplican las ecuaciones diferenciales en otras áreas de las matemáticas y otras ciencias.
    \item Para fomentar el desarrollo de las habilidades investigativas en el estudiantado, se deben implementar proyectos de investigación y si es posible la participación en coloquios o seminarios de investigación.
\end{enumerate}

\textbf{Sugerencias metodológicas:} Adicionalmente, se tienen las siguientes recomendaciones para la persona docente a cargo de este curso:
\begin{enumerate}
    \item Utilizar el recurso tecnológico para reforzar distintos conceptos estudiados en el curso por medio de la visualización, cálculo y experimentación.
    \item Aprovechar momentos adecuados del curso para motivar el estudio por medio de reseñas acerca del desarrollo histórico de las ecuaciones diferenciales y su importancia en aplicaciones dentro y fuera de la matemática. 
    \item En la evaluación, se sugiere utilizar estrategias que promuevan el trabajo constante de parte de las personas estudiantes a lo largo del ciclo lectivo. Por ejemplo, esto se puede hacer por medio de listas de ejercicios recomendados o proyectos.
    \item Dado que hoy en día la mayoría del trabajo en matemática, tanto a nivel académico como profesional, es de naturaleza colaborativa, se sugiere a la persona docente implementar estrategias que promuevan el trabajo en equipo.
    \item Se recomienda el uso de proyectos de investigación bibliográfica para los temas especiales del final del curso. Esto permite un mayor nivel de profundidad además de involucrar más activamente a la persona estudiante en el proceso de aprendizaje.
\end{enumerate}

%\subsubsection*{Guía para las referencias}

\nocite{Cod89}
\nocite{Cod55}
\nocite{Bra93}
\nocite{VE21}
\nocite{Mag23}
\nocite{Per01}
\nocite{Tes01}

\printbibliography[heading=subbibliography,section=\the\value{refsection}]
\end{refsection}
